\documentclass[11pt]{article}
\usepackage[margin=1in]{geometry}
\usepackage{amsmath,amssymb,amsfonts,mathtools,bm}
\usepackage[T1]{fontenc}
\usepackage[utf8]{inputenc}
\title{Aligned Mathematical Structures}
\author{Source-backed grouped formula sample}
\date{}
\begin{document}
\maketitle
\section{Expressions}
\paragraph{Expression 1.} The following expression is taken from a source-backed formula corpus.

\begin{align*}\begin{array}{lll}\delta r_+&=&\delta\left(M+\sqrt{M^2-a^2-Q^2}\right)\\\\&=&\delta M+{1\over2}\,{\left(M^2-a^2-Q^2\right)}^{-{1\over2}}\,2\left(M\,\delta M-a\,\delta a\right)\\\\&=&\delta M+{\left(r_+-M\right)}^{-1}\left(M\,\delta M-a\,\delta a\right)\;,\\\\a\,\delta a&=&a\,\delta\left(J/M\right)={a\over M}\,\delta J-{a^2\over M}\,\delta M\;.\end{array}\end{align*}

\paragraph{Expression 2.} The following expression is taken from a source-backed formula corpus.

\begin{align*}\begin{array}{rl}A\circ_{X}B & = (AX)(X^{ \dagger }B) \\ \\& = (AX^{2/3}X^{1/3})(X^{-1/3}X^{-2/3}B) \\ \\& = X^{1/3}[(X^{-1/3}AX^{2/3})(X^{-2/3}B)] \\ \\& = X^{1/3}[(X^{-1/3}AX^{1/3}X^{1/3})(X^{-1/3}X^{-1/3}B)] \\ \\& = X^{1/3}\{[(X^{-1/3}AX^{1/3})(X^{-1/3}BX^{1/3})]X^{-1/3}\} \\\\ & = X^{1/3}[(X^{-1/3}AX^{1/3})(X^{-1/3}BX^{1/3})]X^{-1/3}. \\\end{array}\end{align*}

\paragraph{Expression 3.} The following expression is taken from a source-backed formula corpus.

\begin{align*}F^{(L)}_{p;r}(q) =q^{-r(r-1)} \times\left\{ \begin{array}{ll}F^{(L)}_{p-2} \left[ \begin{array}{c}{\bf Q}_{r,1} \\ {\bf 0} \end{array} \right]({\bf e}_r | q) &\mbox {if $L \not\equiv r-1$ mod $2$}, \\F^{(L)}_{p-2} \left[ \begin{array}{c}{\bf Q}_{p-r,p} \\ {\bf 0} \end{array} \right]({\bf e}_{p-r} | q) &\text{if $L \equiv r-1$ mod $2$},\end{array}\right.\end{align*}

\paragraph{Expression 4.} The following expression is taken from a source-backed formula corpus.

\begin{align*}\begin{array}{c}u(\tau )=e^{\sqrt{\gamma ^2+m^2/2}\,\tau }\left[ u_0\cosh (\gamma \tau)-\left( u_0\sqrt{\gamma ^2+m^2/2}-\dot u_0\right) \frac 1\gamma \sinh(\gamma \tau )\right] \\v(\tau )=e^{-\sqrt{\gamma ^2+m^2/2}\,\tau }\left[ v_0\cosh (\gamma \tau)+\left( v_0\sqrt{\gamma ^2+m^2/2}+\dot v_0\right) \frac 1\gamma \sinh(\gamma \tau )\right] ,\end{array}\end{align*}

\paragraph{Expression 5.} The following expression is taken from a source-backed formula corpus.

\begin{align*}{\hspace{1mm}^*}\hspace{-1mm}f_{\mu \nu}(x) = - \frac{e_0}{4 \pi \varepsilon_0 c} [ \partial_\mu \vec{n}(x) \wedge \partial_\nu \vec{n}(x) ] \vec{n}(x)=\left( \begin{array}{cccc} 0 & B_x & B_y & B_z \\ -B_x & 0 & \frac{E_z}{c} & \frac{-E_y}{c} \\ -B_y & \frac{-E_z}{c} & 0 & \frac{E_x}{c} \\ -B_z & \frac{E_y}{c} & \frac{-E_x}{c} & 0 \\\end{array} \right)\end{align*}
\end{document}
